\documentclass[11pt, a4paper]{article}
\usepackage[margin=0.8in]{geometry}
\title{Project Title}
\author{Sam \and Charlie \and Anusha \and Mukund}
\date{}

\begin{document}

\maketitle

\section*{Project Goal and Scope}

The primary objective is to develop a photo-realistic 3D graphics rendering system in C++ utilizing advanced Path Tracing Techniques. 

The specific target for the rendering competition is to accurately recreate an image of a chess scene \cite{dartmouth_rendering}. Achieving this level of visual fidelity requires a synthesis of custom geometric modeling, designing material properties, and a Path Tracer that handles shadows, and glass surfaces.


\section*{Accomplishments}
\subsection*{Bidirectional Path Tracing }
We upgraded the core rendering engine to utilize Bidirectional Path Tracing \cite{Vlnas2018}. In traditional Monte Carlo path tracing, rays are exclusively cast from the camera into the scene. Our implementation improves upon this by simultaneously generating sub-paths from both the camera (eye paths) and the emissive light sources (light paths). The length of both sub-paths is determined stochastically, utilizing Russian Roulette for unbiased termination.


\section*{Technical Details}

To evaluate the final pixel color, the integrator attempts to connect the vertices of the camera sub-path with the vertices of the light sub-path. For each attempted connection, a deterministic visibility check (shadow ray) is cast between the two endpoints. If an object occludes the line of sight between the vertices, no color contribution is added for that specific connection. If the path is unoccluded, the radiance is evaluated and accumulated similarly to standard Monte Carlo integration. 

Bi-directional Path-tracing can be defined as 
\begin{equation}
	L_p = \sum_{s=0}^{N_L}\sum_{t=0}^{N_E} w_{s,y} \cdot C_{s,t}
\end{equation}
Where $C_{s,t}$ is the unweighted contribution of a path with $s$ vertices on the light path and $t$ vertices on the eye path. A visibility check between $s^{th}$ point on the eye path and the $t^{th}$ point on the light path is a part of the geometric term in contribution function. And finally, $w_{s,t}$ is the weight function.

Evaluation of $C_{s,t}$ is dependent on the eye and light path; therefore, there are four important cases:
\begin{itemize}
	\item{$s=0,t=0$ , light is directly visible from the eye, the evaluation can be done with direct path between eye the point $x_0$ and $y_0$: $C_{0,0} = L_e(y_0,\vec{y_0x_0}) G(y_0 \leftrightarrow x_0)$}
	\item{$s>0,t=0 , this term is the classic light tracing algorithm$}

\end{itemize}


\section*{Rendering Results and Experiments}


\section*{Discussion}


\bibliographystyle{plain}
\bibliography{references}

\end{document}
